\section{Implications for phenomenology}
\label{sec:pheno}

We now present some initial studies
of the phenomenological implications of NNPDF3.1 PDFs. 
%
Firstly, we summarize the status of PDF uncertainties building
upon the discussion in the previous Section; we discuss the status of 
PDF uncertainties, and then we focus on
the strange and charm PDF. We then discuss PDF luminosities which are
the primary input to hadron collider processes, and predictions for
LHC processes, specifically $W$, $Z$ and Higgs production at
the LHC. 
As elsewhere, only a selection of results is presented here, with
a much larger set is  available online 
as discussed in Section~\ref{sec:delivery}.

\subsection{Improvements in PDF uncertainties}

\label{sec:phenopdfs}

After discussing in Section~\ref{sec:impactnewdata} the impact on
NNPDF3.1 PDFs of each
individual new piece of data and after separating off the effect of the
new methodology, we can now study the combined effect of all the new data 
by comparing PDF uncertainties in NNPDF3.0 and NNPDF3.1. This is done in
Fig.~\ref{fig:pdfunc31vs30}, where we compare relative PDF uncertainties
(all computed with a common normalization) on PDFs in the NNPDF3.0 and
NNPDF3.1 sets, shown as valence (i.e. $q-\bar q$), and sea 
(i.e. $\bar q$) for up and down  and $q+\bar q$ for strange and
charm. Results are shown both for individual PDF flavors, and for the
singlet, valence, and triplet combinations defined in Eq.~(3.4) of
Ref.~\cite{Ball:2016neh}). 

The most visible effect is the very considerable reduction in gluon
uncertainty, which is now at the percent level for almost all $x$. As
discussed in  Section~\ref{sec:impactnewdata}, this is due to the
combination of many mutually consistent constraints on the gluon from
DIS (especially at HERA), $Z$ transverse momentum distributions, jet
production, and top pair production, which taken together cover a
very wide kinematic range. The singlet quark combination, which mixes
with the gluon, shows a comparable improvement for all $x\lsim 0.1$,
but less marked at large $x$.

Interestingly, for quark PDFs the pattern of uncertainties is
different in the flavor basis versus  the ``evolution'' basis as given in Eq.~(3.4) of
Ref.~\cite{Ball:2016neh}). Specifically, 
the aforementioned reduction in uncertainty on the singlet combination
is not seen in any of the 
light quark valence distributions, which generally have comparable
uncertainties in NNPDF3.1 and NNPDF3.0. This is due to the fact that
flavor separation is somewhat more uncertain in NNPDF3.1, due to the
fact that charm is now independently parametrized. This is compensated
by the availability of more experimental information (in particular
LHCb and ATLAS data), but not at small and very large $x$. Indeed, the
valence and triplet distributions  have generally somewhat larger
uncertainties  in
NNPDF3.1 than in NNPDF3.0, except for $x\sim0.1$.

%%%%%%%%%%%%%%%%%%%%%%%%%%%%%%%%%%%%%%%%%%%%%%%%%%%%%%%%%%%%%%%%%%%%%
\begin{figure}[t!]
\begin{center}
  \includegraphics[scale=0.38]{images/plots/valence_nnpdf30nnlo.pdf}
  \includegraphics[scale=0.38]{images/plots/valence_nnpdf31nnlo.pdf}
  \includegraphics[scale=0.38]{images/plots/sea_nnpdf30nnlo.pdf}
  \includegraphics[scale=0.38]{images/plots/sea_nnpdf31nnlo.pdf}
  \includegraphics[scale=0.38]{images/plots/evolbasis3_nnpdf30nnlo.pdf}
  \includegraphics[scale=0.38]{images/plots/evolbasis3_nnpdf31nnlo.pdf}
  \caption{\small 
Comparison of relative uncertainties on NNPDF3.0 (left) and NNPDF3.1
  (right) NNLO PDFs, normalized to the NNPDF3.1
  NNLO central value. The two light quark valence PDFs and the gluon
   are shown  (top) along with all individual sea PDFs (center) and the
   singlet, valence and isospin triplet combinations (bottom).
    %
\label{fig:pdfunc31vs30}}
\end{center}
\end{figure}
%%%%%%%%%%%%%%%%%%%%%%%%%%%%%%%%%%%%%%%%%%%%%%%%%%%%%%%%%%%%%%%%%%%%%%

This  pattern is also seen in  sea uncertainties. At small
$x\lsim 10^{-2}$ they are all comparable, and all (except for strangeness) 
rather smaller than
the corresponding uncertainties in NNPDF3.0, since they are driven by
the mixing of the singlet and the gluon through perturbative
evolution at small $x$~\cite{das}. Even the charm PDF, now
independently parametrized, has
a smaller uncertainty. In the large-$x$ region, instead, the less
accurate knowledge of flavor separation kicks in, and  relative uncertainties 
are larger.
Clearly, in NNPDF3.0
charm had an unnaturally small uncertainty, since at large $x$ 
perturbatively-generated charm is tied to the gluon. In NNPDF3.1, instead, the
hierarchy of uncertainties on sea PDFs in the valence region is what
one would expect, with up and down known most accurately, and strange
and charm affected by increasingly large uncertainties. The
uncertainty in NNPDF3.1 on the up and especially down sea components is a little 
increased, but 
still comparable to NNPDF3.0, while the uncertainty in strangeness is
stable and that on 
charm significantly increased, as it should be given that large-$x$
charm is largely unconstrained by data. In this respect, it is
interesting to observe that a more accurate determination of charm
and other sea PDFs at large $x$ can be achieved through the inclusion
of the EMC dataset as discussed in Section~\ref{sec:emc} above. Future LHC
data on processes such as $Z+c$ production may confirm the reliability
of the EMC dataset.

\subsection{The strange PDF}
\label{sec:strangeness}

Whereas there is broad consensus on the size, i.e. the central value,
of up and down PDFs, for which there is good agreement between existing
determinations within their small uncertainty, the size of the strange
PDF has been the object of some controversy, which we revisit here in
view of NNPDF3.1 results.
Specifically, the strange fraction of proton quark sea, defined as
\be
\label{eq:rs}
R_s (x,Q^2)=\frac{ s(x,Q^2)+\bar{s}(x,Q^2) }{ \bar{u}(x,Q^2)+\bar{d}(x,Q^2) } \, .
\ee
and the corresponding ratio of momentum fractions
\be
\label{eq:rsint}
K_s (Q^2)=\frac{\int_0^1 dx\, x \lp s(x,Q^2)+\bar{s}(x,Q^2)\rp }{ 
\int_0^1 dx\, x\lp \bar{u}(x,Q^2) + \bar{d}(x,Q^2)\rp }  \, ,
\ee
have been traditionally assumed to be significantly smaller than one,
and in PDF sets produced before the strange PDF could be extracted
from the data, such as e.g. NNPDF1.0~\cite{Ball:2008by}, it was often
assumed that $R_s\sim\frac{1}{2}$, for all $x$, and thus also 
$K_s\sim\frac{1}{2}$. This level of strangeness suppression is indeed found in
many recent global PDF sets, in which the strongest handle on the strange
PDF is provided by deep-inelastic neutrino inclusive $F_2$ and  charm
$F_2^c$ (``dimuon'') data. 

This was challenged in Ref.~\cite{Aad:2012sb} where, on the
basis of ATLAS $W$ and $Z$ production data, combined with HERA DIS
data, 
it was argued instead that,
in the measured region, the strange fraction $R_s$ is of order one. In
Refs.~\cite{Ball:2012cx,Ball:2015oha}, respectively based on the
NNPDF2.3, and NNPDF3.0 global analyses, both of which
includes the data of Ref.~\cite{Aad:2012sb}, it was concluded that
whereas the ATLAS data do favor a larger strange PDF, they have a moderate
impact on the global PDF determination 
due to large uncertainties, and also, that if
the strange PDF is only determined from HERA and ATLAS data, the
central value is consistent with the conclusion of
Ref.~\cite{Aad:2012sb}, but the uncertainty is large enough to lead to
agreement with the suppressed strangeness of the global PDF sets to within
one sigma. In Ref.~\cite{Ball:2015oha} it was also shown that the CMS
$W+c$ production data~\cite{Chatrchyan:2013uja}, which were included
there for the first time and which are also included in NNPDF3.1, 
though only in the NLO
determination because of lack of knowledge of the NNLO corrections,
have a negligible impact due to their large uncertainties. 

As we discussed in Section~\ref{sec:atlaswz},  ATLAS $W$ and $Z$
production data have been supplemented by the rather more accurate
dataset of Ref.~\cite{Aaboud:2016btc}, also claimed to 
favor enhanced strangeness. Indeed, we have seen in
Section~\ref{sec:atlaswz} that strangeness is significantly enhanced by
the inclusion of these data, and also, in Section~\ref{sec:results-mc},
that this enhancement  can be accommodated in the global PDF
determination 
thanks to the independently parametrized  charm PDF, which is a new
feature to NNPDF3.1.
It is thus interesting to re-asses strangeness in NNPDF3.1, by
comparing theoretically motivated choices of dataset: we will thus compare
to the
previous NNPDF3.0 results for strangeness obtained using the default
NNPDF3.1, the 
collider-only PDF set of Section~\ref{sec:collideronly}, which can be
considered to be theoretically more reliable, and a PDF set which we
have constructed by using NNPDF3.1 methodology, but only including all
HERA inclusive structure function data from
Tab.~\ref{tab:completedataset} and the ATLAS data of
Ref.~\cite{Aaboud:2016btc}. Because inclusive DIS data alone cannot
determine separately strangeness~\cite{Forte:2010dt} this is then a
determination of strangeness which fully relies on the ATLAS data.

In Table~\ref{tab:rs}  we show NNLO
results, obtained using these different PDF sets,
 for $R_s(x,Q)$ Eq.~\ref{eq:rs} at
 $Q=1.38~{\rm GeV}$ (thus below charm threshold) and $Q=m_Z$ and
$x=0.023$, an $x$ value chosen by ATLAS in order to maximize
sensitivity. Results are also compared to that of
Ref.~\cite{Aaboud:2016btc}.
A graphical representation of the table is in Fig.~\ref{fig:xRsplot}.

%%%%%%%%%%%%%%%%%%%%%%%%%%%%%%%%%%%%%%%%%%%%%%%%%%%%%%%%%%%%%%%%%%%
\begin{table}[t]
  \begin{center}
    \small
        \renewcommand{\arraystretch}{1.2}
\begin{tabular}{|l|c|c|}
\hline
\centering PDF set &  $R_s(0.023,1.38~\mathrm{GeV})$ & $R_s(0.023,M_{Z})$ \\
\hline
\hline
NNPDF3.0    & 0.45$\pm$0.09    & 0.71$\pm$0.04 \\
NNPDF3.1     & 0.59$\pm$0.12  & 0.77$\pm$0.05 \\
NNPDF3.1 collider-only   & 0.82$\pm$0.18  & 0.92$\pm$0.09 \\
NNPDF3.1 HERA + ATLAS $W,Z$   & 1.03$\pm$0.38  & 1.05$\pm$0.240\\	
\hline
{\tt xFitter} HERA + ATLAS $W,Z$ (Ref.~\cite{Aaboud:2016btc}) 	& $1.13
\,{}^{+0.11}_{-0.11}$& - \\
\hline
\end{tabular}
\end{center}
\caption{\small \label{tab:rs} 
  The strangeness fraction $R_s(x,Q)$  Eq.~(\ref{eq:rs}) at $x=0.023$,
  at a low scale and a high scale. We show results obtained using  
NNPDF3.0, and NNPDF3.1 baseline, collider-only and  HERA+ATLAS $W,Z$
sets, compared to the {\tt xFitter}  ATLAS value  Ref.~\cite{Aaboud:2016btc}.
}
\end{table}
%%%%%%%%%%%%%%%%%%%%%%%%%%%%%%%%%%%%%%%%%%%%%%%%%%%%%%%%%%%%%%%%%%%

%%%%%%%%%%%%%%%%%%%%%%%%%%%%%%%%%%%%%%%%%%%%%%%%%%%%%%%%%%%%%%%%%%%%%
\begin{figure}[t]
\begin{center}
  \includegraphics[scale=0.35]{images/plots/Rsplot.pdf}
  \includegraphics[scale=0.35]{images/plots/Rsplot-highQ.pdf}
  \caption{\small Graphical representation of the results
    of Table~\ref{tab:rs}.
\label{fig:xRsplot}}
\end{center}
\end{figure}
%%%%%%%%%%%%%%%%%%%%%%%%%%%%%%%%%%%%%%%%%%%%%%%%%%%%%%%%%%%%%%%%%%%%%%

First, comparison of  the NNPDF3.1 HERA+ATLAS $W,Z$ result with that of
Ref.~\cite{Aaboud:2016btc}, based on the same data, shows agreement at
the one-sigma level, with a similar central value and a
greatly increased 
uncertainty, about four times larger, most likely because of the more
flexible parametrization and because of independently parametrizing charm.
Second, strangeness in NNPDF3.1 is rather larger than in NNPDF3.0: as
we have shown in Sects.~\ref{sec:results-mc},\ref{sec:atlaswz} this is
largely due to the effect of the
ATLAS~$W,Z$~2011 data, combined with determining  charm from the data:
indeed, it is clear
from Fig.~\ref{fig:31-nnlo-old-vs-new} that the new data and new
methodology both lead to strange enhancement, with the former effect
dominant but the latter not negligilbe. 
This enhancement is more marked in the collider-only PDF
set, which leads to a value which is very close to that coming from
the ATLAS data. This suggests some tension between 
strangeness preferred by collider data and the rest of the dataset,
i.e., most likely, neutrino data.


It is interesting to repeat this analysis for the full $x$ range. 
This is done in Fig.~\ref{fig:xRs-abs-atlaswz2011}, where
$R_s(x,Q)$ Eq.~(\ref{eq:rs}) is plotted as a function
of $x$ again at low and high scales, now only including NNPDF3.0, and
the default and collider-only versions of NNPDF3.1. It is clear that
in the collider-only PDF set strangeness is largely unconstrained   at
large $x$, whereas the global fit is constrained by neutrino data to
have a suppressed value $R_s\sim 0.5$. At lower $x$ we see the tension
between this and the constraint
from the collider data, which prefer a larger value.

%%%%%%%%%%%%%%%%%%%%%%%%%%%%%%%%%%%%%%%%%%%%%%%%%%%%%%%%%%%%%%%%%%%%%
\begin{figure}[t]
\begin{center}
  \includegraphics[scale=0.37]{images/plots/xRs-abs-atlaswz2011lowQ.pdf}
  \includegraphics[scale=0.37]{images/plots/xRs-abs-atlaswz2011highQ.pdf}
   \includegraphics[scale=0.37]{images/plots/xRs-abs-atlaswz2011lowQ-global.pdf}
  \includegraphics[scale=0.37]{images/plots/xRs-abs-atlaswz2011highQ-global.pdf}
  \caption{\small The strangeness ratio $R_s(x,Q)$ Eq.~(\ref{eq:rs})
    as a function of 
    $x$ for two values of $Q$, $Q=1.38$ GeV (left)
    and $Q=m_Z$ (right).
    %
     Results are shown comparing NNPDF3.1 to  NNPDF3.1 and the
     collider-only NNPDF3.1 (top), and to CT14 and MMHT (bottom).
\label{fig:xRs-abs-atlaswz2011}}
\end{center}
\end{figure}
%%%%%%%%%%%%%%%%%%%%%%%%%%%%%%%%%%%%%%%%%%%%%%%%%%%%%%%%%%%%%%%%%%%%%%

In Fig.~\ref{fig:xRs-abs-atlaswz2011} we also compare the strangeness ratio $R_s(x,Q)$ of
NNPDF3.1 with that of CT14 and MMHT14.
%
We find that there is good consistency in the entire range of $x$,
while the PDF errors
in NNPDF3.1 are typically smaller than those of the other two sets,
especially at  large scales.
It is also interesting to note how in NNPDF3.1 the PDF uncertainties
in the ratio $R_s$ blow up at very large $x$, reflecting the lack
of direct information on strangeness in that kinematic region.


%%%%%%%%%%%%%%%%%%%%%%%%%%%%%%%%%%%%%%%%%%%%%%%%%%%%%%%%%%%%%%
\begin{table}[t]
  \begin{center}
        \renewcommand{\arraystretch}{1.2}
\begin{tabular}{|l|c|c|}
\hline
\centering PDF set &  $K_s(Q=1.38~\mathrm{GeV})$ & $K_s(Q=M_{\rm Z})$ \\
\hline
\hline
NNPDF3.0   &  $0.45 \pm 0.07$  & $0.72 \pm 0.04$ \\
NNPDF3.1    & $0.53 \pm 0.07$  &  $ 0.75 \pm 0.04 $  \\
NNPDF3.1 collider-only   & $3.4 \pm 2.5$  & $1.5 \pm 0.6$  \\
NNPDF3.1 HERA + ATLAS $W,Z$   & $-1.0 \pm 7.0$  & $2.8 \pm 1.7$ \\
\hline
\end{tabular}
\end{center}
\caption{\small \label{tab:rsint} 
  The strangeness momentum fraction Eq.~(\ref{eq:rsint}) 
  at a low scale and a high scale. We show results obtained using  
NNPDF3.0, and NNPDF3.1 baseline, collider-only and  HERA+ATLAS $W,Z$
PDF sets.}
\end{table}
%%%%%%%%%%%%%%%%%%%%%%%%%%%%%%%%%%%%%%%%%%%%%%%%%%%%%%%%%%%%%%%%%%


We now turn to  the strange momentum fraction $K_s(Q^2)$
Eq.~(\ref{eq:rsint}); values for the same PDF sets and scales are
shown in Table~\ref{tab:rsint}. Results are quite similar to those
found from the analysis of  Table~\ref{tab:rs}.
For the NNPDF3.1 collider-only and especially the HERA + ATLAS $W,Z$ fits, the
central value of $K_s$ is  
unphysical, with a huge uncertainty; essentially, all one can say is
that the strange momentum fraction $K_s$ is completely uncertain. 
This shows rather dramatically that the relatively precise values 
in Table~\ref{tab:rs} only hold in a rather narrow $x$ range.
It will be interesting to see whether more LHC data, possibly leading
to a competitive collider-only fit,  will confirm
strangeness enhancement and allow for an accurate determination of
strangeness in a wider range of $x$.

\subsection{The charm content of the proton revisited}
\label{sec:phenocharm}

The charm content of the proton determined by fitting the charm PDF
was quantified within the
NNPDF global analysis framework in~\cite{Ball:2016neh}, where results
obtained when  charm is independently parametrized, or perturbatively
generated , were compared for the first time.
%
The analysis there was performed at NLO only, and the dataset was very similar to that
of the NNPDF3.0 fit.
%
We now re-examine the fitted charm PDF at NNLO in perturbative QCD, and in the 
context of the inclusion of the new datasets,
in particular top, and LHCb and ATLAS electroweak boson production,
which sizably affects and
constrain the charm PDF. 

Indeed, we have seen in
Section~\ref{sec:disentangling}, in particular
Fig.~\ref{fig:31-nnlo-old-vs-new}, that the new data added in
NNPDF3.1 considerably reduces the charm PDF uncertainty, but also affects its
central value, which changes by more than one sigma at large $x$.
Also, in Ref.~\cite{Ball:2016neh} charm at large $x$ was mostly
constrained by the EMC data which we discussed in Section~\ref{sec:emc}
and which are not included in the default NNPDF3.1 PDF
determination. As we have seen in Section~\ref{sec:emc} these data still
have a significant impact on charm, hence a re-assessment of charm
determination is in order both when this dataset is included and when it
is not. We therefore now compare results obtained using the
default NNPDF3.1 NNLO set, the modified version of that in which charm
is generated perturbatively as discussed in
Section~\ref{sec:results-mc}, and the modified set in which the EMC data are added
to the NNPDF3.1 dataset, as discussed in Section~\ref{sec:emc}.

In Table~\ref{tab:momsr} we show the charm momentum fraction, defined as
\be
\label{eq:charmmomfrac}
C(Q^2) \equiv \int_0^1\, dx\,\lp xc(x,Q^2)+x\bar{c}(x,Q^2)\rp \ ,
\ee
for two values of $Q^2$, at the charm threshold $Q=m_c$, and at
$Q=m_Z$, computed from using these PDF sets.
A graphical representation of the results from 
Table~\ref{tab:momsr} is shown in Fig.~\ref{fig:xCplot}.
There is a very large difference, by two orders of
magnitude, between the uncertainty on the momentum fraction, according
to whether charm is independently parametrized, or perturbatively generated.
However, the central values
agree with each other
within the large uncertainty determination when charm is parametrized, with
the corresponding central value only slightly larger than that when
charm is  perturbative
(though hugely different on the scale of the uncertainty on the
perturbatively generated result).
Upon adding the EMC data the
uncertainty is reduced by about a factor of three, and the central
value somewhat increased, consistently with the effect of this dataset on
the charm PDF discussed in Section~\ref{sec:emc}.

We can interpret the  difference between the total momentum fraction
when charm is independently parametrized and determined from the data, 
and that when charm is perturbatively generated, as the
``intrinsic'' (i.e. non-perturbative) charm momentum fraction.
Including EMC data when charm is parametrized,
at  $Q=m_c$  we find that it is given by
$C(m_c)_{\rm FC}-C(m_c)_{\rm PC} =  \left(0.16 \pm 0.14\right)$\%: this provides evidence
for a small intrinsic charm component in the proton at the
one $\sigma$ level, somewhat improving the estimates of
Ref.~\cite{Ball:2016neh}, but 
still with a considerable degree of uncertainty.
%
Our estimate for the non-perturbative charm component of the proton
is considerably smaller than those allowed in the CT14IC model analysis~\cite{Dulat:2013hea}, 
also shown in Fig.~\ref{fig:xCplot}. However, it is larger than expected 
in Ref.~\cite{Jimenez-Delgado:2014zga}, where an upper bound of $0.5$\% at the 
four-$\sigma$ level is claimed. Both these analyses have difficulty fitting the EMC charm 
structure function data, due to an overly restrictive functional form
for the charm PDF. At high scale all estimates of $C(Q)$ are dominated by the perturbative
growth of the charm PDF at small $x$, but the one-$\sigma$ excess in the fit 
with EMC data persists, though as an ever diminishing fraction of the whole. 
This is demonstrated very clearly in Fig.~\ref{fig:momfrac}, where we plot the dependence 
of the charm momentum fraction $C(Q^2)$ on the scale $Q$. 


%%%%%%%%%%%%%%%%%%%%%%%%%%%%%%%%%%%%%%%%%%%%%%%%%%%%%%%%%%%%%%%%%%%%%%%
\begin{table}[t]
  \center
   \renewcommand{\arraystretch}{1.4}
  \begin{tabular}{|l|c|c|}
    \hline
    PDF set & $C(m_c)$ & $C(m_Z)$ \\
    \hline
    \hline
    NNPDF3.1           & $\lp 0.26 \pm 0.42\rp $\%  &  $\lp 3.8 \pm 0.3\rp $\%\\
    NNPDF3.1 with pert. charm        & $\lp 0.176 \pm 0.005\rp $\%  &  $\lp 3.73 \pm 0.02\rp $\% \\
    NNPDF3.1 with EMC data  &   $\lp 0.34 \pm 0.14\rp $\% & $\lp 3.8  \pm 0.1\rp $\% \\
          \hline
  \end{tabular}
  \caption{\small \label{tab:momsr}
    The charm momentum fraction  $C(Q^2)$, Eq.~(\ref{eq:charmmomfrac}),
    just above the charm threshold $Q=m_c$~GeV and at
    $Q=m_Z$. Results are shown for NNPDF3.1, and its modified versions
    in which EMC data are added to the dataset,  or charm is not fitted. 
    %
  }
  \end{table}
%%%%%%%%%%%%%%%%%%%%%%%%%%%%%%%%%%%%%%%%%%%%%%%%%%%%%%%%%%%%%%%%%%%%%%%

%%%%%%%%%%%%%%%%%%%%%%%%%%%%%%%%%%%%%%%%%%%%%%%%%%%%%%%%%%%%%%%%%%%%%
\begin{figure}[t]
\begin{center}
  \includegraphics[scale=0.35]{images/plots/Cplot.pdf}
  \includegraphics[scale=0.35]{images/plots/Cplot-highQ.pdf}
  \caption{\small Graphical representation of the results
    for $C(Q^2)$ from Table~\ref{tab:momsr} 
    and $Q=m_c$ GeV (left) and $Q=m_Z$ (right). Model estimates from
    Ref.~\cite{Dulat:2013hea} are also shown.
\label{fig:xCplot}}
\end{center}
\end{figure}
%%%%%%%%%%%%%%%%%%%%%%%%%%%%%%%%%%%%%%%%%%%%%%%%%%%%%%%%%%%%%%%%%%%%%%



%%%%%%%%%%%%%%%%%%%%%%%%%%%%%%%%%%%%%%%%%%%%%%%%%%%%%%%%%%%%%%%%%%%%%
\begin{figure}[t]
\begin{center}
  \includegraphics[scale=0.30]{images/plots/momfrac.pdf}
  \caption{\small The charm momentum fraction of Tab.~\ref{tab:momsr}
    plotted  as a function of  scale $Q$.
\label{fig:momfrac}}
\end{center}
\end{figure}
%%%%%%%%%%%%%%%%%%%%%%%%%%%%%%%%%%%%%%%%%%%%%%%%%%%%%%%%%%%%%%%%%%%%%%

The origin of these values of the charm momentum fraction can be
understood by directly  comparing the charm PDF, which is done
in Fig.~\ref{fig:PDFlargexcharm}, again 
 just above threshold $Q=m_c$ and at $Q=m_Z$, in the latter case
as a ratio to the baseline NNPDF3.1 result.
The agreement of the charm  momentum fraction when it is
perturbatively generated or when it is parametrized and determined
from the data is related to the fact that, when parametrized, 
the best-fit charm
has qualitatively the same shape as charm generated perturbatively at
NNLO, as observed in Ref.~\cite{Ball:2016neh} and
Section~\ref{sec:results-mc} above. 
However, at threshold $Q=m_c$ GeV the best-fit charm
is larger than the perturbative component at large $x$, $x\gsim 0.2$,
albeit with large uncertainties, that are somewhat reduced when the EMC
dataset is added, without a significant change in shape. 
Upon addition of the EMC data, in the
medium-small  $10^{-2}\lsim x\lsim 10^{-1}$ the PDF is pushed at the upper edge
of the uncertainty band before addition, with considerably reduced
uncertainty. The unrealistically small uncertainty on the
perturbatively generated charm PDF is apparent, and also the reduction
in uncertainty due to the EMC data for all $x\gsim 10^{-3}$, already
discussed in Section~\ref{sec:emc}.

%%%%%%%%%%%%%%%%%%%%%%%%%%%%%%%%%%%%%%%%%%%%%%%%%%%%%%%%%%%%%%%%%%%%%
\begin{figure}[t]
  \begin{center}
    \includegraphics[scale=0.35]{images/plots/xc-abs-charmcomp-lowQ.pdf}
  \includegraphics[scale=0.35]{images/plots/xc-charmcomp-highQ.pdf}
  \includegraphics[scale=0.35]{images/plots/xc-ERR-charmcomp-lowQ.pdf}
  \includegraphics[scale=0.35]{images/plots/xc-ERR-charmcomp-highQ.pdf}
    \caption{\small Comparison of the charm PDF at the scale and for
      the PDF sets of Tab.~\ref{tab:momsr}. Both the PDF (top) and the
      relative uncertainty (bottom) are shown.
\label{fig:PDFlargexcharm}}
\end{center}
\end{figure}
%%%%%%%%%%%%%%%%%%%%%%%%%%%%%%%%%%%%%%%%%%%%%%%%%%%%%%%%%%%%%%%%%%%%%%

The origin of the differences between the charm PDF when
perturbatively generated, or parametrized
and determined from the data,
and a possible decomposition of the latter into an ``intrinsic''
(non-perturbative) and a perturbative component can be understood by
studying their scale dependence close to threshold, in analogy to a
similar analysis presented in Ref.~\cite{Ball:2016neh}.
This is done 
in Fig.~\ref{fig:charmPDFnf4}, where the charm PDF is shown
(in the $n_f=4$ scheme) as a function of $x$ both when charm is
parametrized and when it is perturbatively generated.
On the one hand, the large $x\gsim0.1$ component of the
 charm PDF is essentially scale independent: perturbative
charm vanishes identically in this region, so the fitted result in
this region may
    be interpreted as being of non-perturbative origin, i.e. ``intrinsic''.

On the other hand, for smaller $x$ the charm PDF depends
strongly on scale. When perturbatively generated, it  is sizable and
positive already 
at $Q\sim 2$~GeV for all $x\gsim 10^{-3}$, while at threshold it 
becomes large and negative for all $x\lsim10^{-2}$, possibly because
of large unresummed small-$x$ contributions. The best-fit parametrized
 charm PDF, within its larger uncertainty, is rather flatter and 
smaller in modulus  
essentially for all $x\lsim 10^{-1}$, 
except at the scale-dependent point at which
perturbative charm changes sign. This difference in shape between
fitted charm and perturbative charm for all $x\lsim 0.1$ is rather
larger than the uncertainty on either  perturbative or  fitted charm. 
This observation seems to support the conclusion
that the assumption that charm is perturbatively generated might be a
source of bias.

%%%%%%%%%%%%%%%%%%%%%%%%%%%%%%%%%%%%%%%%%%%%%%%%%%%%%%%%%%%%%%%%%%%%%
\begin{figure}[t]
\begin{center}
  \includegraphics[scale=0.37]{images/plots/charm-nf4-Qdep-log.pdf}
  \includegraphics[scale=0.37]{images/plots/charm-nf4-Qdep-lin.pdf}
  \includegraphics[scale=0.37]{images/plots/charm-nf4-Qdep-log-nFC.pdf}
  \includegraphics[scale=0.37]{images/plots/charm-nf4-Qdep-lin-nFC.pdf}
  \caption{\small The charm PDF in the $n_f=4$ scheme at small $x$
    (left) and large $x$ (right plot) for different values of $Q$,
    in the NNPDF3.1 NNLO PDF set (top) and when assuming that charm is
    perturbatively generated (bottom).
\label{fig:charmPDFnf4}}
\end{center}
\end{figure}
%%%%%%%%%%%%%%%%%%%%%%%%%%%%%%%%%%%%%%%%%%%%%%%%%%%%%%%%%%%%%%%%%%%%%%

The updated NNPDF3.1 analysis is consistent with the conclusion of our 
earlier charm studies~\cite{Ball:2016neh}, namely 
that a non-perturbative charm component in the proton is consistent with
global data, though the new high-precision collider data that we include now
sets more stringent bounds on how large this component can be.
%
In particular, once the EMC charm structure function dataset is added,
we see from Table~\ref{tab:momsr} that a non-perturbative charm
momentum fraction of $\simeq 0.5\%$ represents a 
deviation from our best fit value of around two to three sigma. Thus models of 
intrinsic charm which carry as much as 1\% of the proton's momentum 
are strongly disfavored by currently available data.

\subsection{Parton luminosities}
\label{sec:lumis}

After analytically discussing the phenomenological implication of
individual PDFs and their uncertainties we now turn to parton
luminosities (defined as in Ref.~\cite{Mangano:2016jyj}) which drive
hadron collider processes. 
Parton luminosities from the NNPDF3.0 and NNPDF3.1 NNLO sets for the
LHC~$13$~TeV are compared
in   Fig.~\ref{fig:lumi-31-vs-30}, and their
uncertainties are displayed in  Fig.~\ref{fig:lumi-nnpdf31-2d} as a
two-dimensional contour plot as a function of the invariant mass $M_y$
and rapidity $y$ of the final state, all normalized to the NNPDF3.1
central value. We show results for the quark-quark, quark-antiquark,
  gluon-gluon  and gluon-quark luminosities, relevant for
the measurement of final states which do not couple to individual
flavors (such as $Z$ or Higgs). In the uncertainty plot we also show
for reference the up-antidown luminosity, relevant e.g. for $W^+$
production. 

%%%%%%%%%%%%%%%%%%%%%%%%%%%%%%%%%%%%%%%%%%%%%%%%%%%%%%%%%%%%%%%%%%%%%
\begin{figure}[t]
\begin{center}
  \includegraphics[scale=0.38]{images/plots/q2_31-vs-30.pdf}
  \includegraphics[scale=0.38]{images/plots/qq_31-vs-30.pdf}
  \includegraphics[scale=0.38]{images/plots/gg_31-vs-30.pdf}
   \includegraphics[scale=0.38]{images/plots/qg_31-vs-30.pdf}
   \caption{\small Comparison of parton luminosities with the NNPDF3.0
     and NNPDF3.1 NNLO PDF sets for the LHC~$13$ TeV. From left to
     right and from top to bottom 
 quark-antiquark, quark-quark, gluon-gluon
     and quark-gluon PDF luminosities are shown.     %
     Results are shown normalized to the central value of NNPDF3.1.
\label{fig:lumi-31-vs-30}}
\end{center}
\end{figure}
%%%%%%%%%%%%%%%%%%%%%%%%%%%%%%%%%%%%%%%%%%%%%%%%%%%%%%%%%%%%%%%%%%%%%%


%%%%%%%%%%%%%%%%%%%%%%%%%%%%%%%%%%%%%%%%%%%%%%%%%%%%%%%%%%%%%%%%%%%%%
\begin{figure}[h]
\begin{center}
   \includegraphics[scale=0.40]{images/plots/qqbar_plot_lumi2d_uncertainty_30.pdf}   
   \includegraphics[scale=0.40]{images/plots/qqbar_plot_lumi2d_uncertainty.pdf}   
\includegraphics[scale=0.40]{images/plots/qq_plot_lumi2d_uncertainty_30.pdf}
\includegraphics[scale=0.40]{images/plots/qq_plot_lumi2d_uncertainty.pdf}
\includegraphics[scale=0.40]{images/plots/gg_plot_lumi2d_uncertainty_30.pdf}
\includegraphics[scale=0.40]{images/plots/gg_plot_lumi2d_uncertainty.pdf}
  \includegraphics[scale=0.40]{images/plots/gq_plot_lumi2d_uncertainty_30.pdf}
  \includegraphics[scale=0.40]{images/plots/gq_plot_lumi2d_uncertainty.pdf}
 \includegraphics[scale=0.40]{images/plots/udbar_plot_lumi2d_uncertainty_30.pdf}
 \includegraphics[scale=0.40]{images/plots/udbar_plot_lumi2d_uncertainty_31.pdf}
   \caption{\small The relative uncertainty on the luminosities of
     Fig.~\ref{fig:lumi-31-vs-30}, plotted as  a function of
     the invariant mass $M_X$ and the rapidity $y$ of the final
     state; the left plots show results for NNPDF3.0 and the right
     plots for NNPDF3.1 (upper four rows). The bottom row shows
     results for the up-antidown luminosity.
\label{fig:lumi-nnpdf31-2d}}
\end{center}
\end{figure}
%%%%%%%%%%%%%%%%%%%%%%%%%%%%%%%%%%%%%%%%%%%%%%%%%%%%%%%%%%%%%%%%%%%%%%


Two features of this comparison are apparent. First, quark
luminosities are generally larger for all invariant masses, while the
gluon luminosity is somewhat enhanced for smaller invariant and
somewhat suppressed for  larger invariant masses in
NNPDF3.1 in comparison to NNPDF3.0.
The size of the shift in the quark sector is of order of one sigma or
sometimes even larger, while for the gluon is generally rather below
the one-sigma level. This of course reflects the pattern seen in
Section~\ref{sec:PDFcomparisons} for PDFs,
see in particular Fig.~\ref{fig:31-nnlo-vs30}.
Secondly, uncertainties are greatly reduced in NNPDF3.1 in comparison to
NNPDF3.0. This reduction is impressive and apparent in the plots of
Fig.~\ref{fig:lumi-nnpdf31-2d}, where it is clear that while
uncertainties were typically of order 5\% in most of phase space for
NNPDF3.0, they are now of the order of 1-2\% in a wide central
rapidities range
$|y|\lsim 2$ and for final state masses $100$~GeV$\lsim M_x\lsim 1$~TeV.
This is a direct consequence of  the reduction in uncertainties on both the
gluon and quark singlet PDF discussed in Section~\ref{sec:phenopdfs}
above. Indeed, luminosities which are sensitive to the flavor
decomposition, such as the up-antidown luminosity, also shown in
Fig.~\ref{fig:lumi-nnpdf31-2d}, do not display a significant reduction
in uncertainties when going to NNPDF3.0 to NNPDF3.1.


We next compare  NNPDF3.1 with CT14 and MMHT14: results are 
shown in Fig.~\ref{fig:lumi-31-vs-global}.
For the quark-quark luminosities, we find good agreement, while
for quark-antiquark  there is a somewhat bigger spread in central
values though still  agreement at the one-sigma level.
Agreement becomes marginal at large masses, $M_X\gsim 2$ TeV, reflecting
the limited knowledge of the large-$x$  PDFs. For the 
gluon-gluon and gluon-quark  channels
we find reasonable agreement for masses up to $M_X\simeq 600$ GeV, relevant for 
precision physics at the LHC, but rather worse agreement for larger
masses, relevant for BSM searches,
in particular between NNPDF3.1 and MMHT14.
%
Of course it should be kept in mind that NNPDF3.1 has a wider
dataset and a larger number of independently parametrized PDFs than
MMHT14 and CT14, hence the situation may change in the future once
all global PDF sets are updated.

Next, in Fig.~\ref{fig:lumi-31-vs-abmp} we compare to
ABMP16 PDFs. In this case, we show results corresponding both to the
default ABMP16 set, which has  $\alpha_s(m_Z)=0.1147$, and to the set with
the common  $\alpha_s(m_Z)=0.118$ adopted so far in all comparison.
     %
     While there
     are sizable differences between NNPDF3.1 and ABMP16 when
     the default ABMP16 value  $\alpha_s(m_Z)=0.1147$ is used, especially for 
     the gluon-gluon luminosity, the agreement improves when
     $\alpha_s(m_Z)=0.118$ is adopted also for ABMP16. 
   However, ABMP16 luminosities have very small uncertainties
  at low and high $M_X$, presumably a consequence of an
  over-constrained parametrization, and of using a Hessian approach
  but with no tolerance, as discussed in Section~\ref{sec:PDFcomparisons}.

\clearpage

%%%%%%%%%%%%%%%%%%%%%%%%%%%%%%%%%%%%%%%%%%%%%%%%%%%%%%%%%%%%%%%%%%%%%
\begin{figure}[t]
\begin{center}
  \includegraphics[scale=0.38]{images/plots/q2_31-vs-global.pdf}
  \includegraphics[scale=0.38]{images/plots/qq_31-vs-global.pdf}
  \includegraphics[scale=0.38]{images/plots/gg_31-vs-global.pdf}
   \includegraphics[scale=0.38]{images/plots/qg_31-vs-global.pdf}
   \caption{\small Same as Fig.~\ref{fig:lumi-31-vs-30}, now comparing
     NNPDF3.1 NNLO  to CT14 and MMHT14.
\label{fig:lumi-31-vs-global}}
\end{center}
\end{figure}
%%%%%%%%%%%%%%%%%%%%%%%%%%%%%%%%%%%%%%%%%%%%%%%%%%%%%%%%%%%%%%%%%%%%%%

%%%%%%%%%%%%%%%%%%%%%%%%%%%%%%%%%%%%%%%%%%%%%%%%%%%%%%%%%%%%%%%%%%%%%
\begin{figure}[t]
\begin{center}
  \includegraphics[scale=0.35]{images/plots/q2_31-vs-abmp.pdf}
  \includegraphics[scale=0.35]{images/plots/qq_31-vs-abmp.pdf}
  \includegraphics[scale=0.35]{images/plots/gg_31-vs-abmp.pdf}
   \includegraphics[scale=0.35]{images/plots/qg_31-vs-abmp.pdf}
   \caption{\small Same as Fig.~\ref{fig:lumi-31-vs-global}, now
     comparing to  ABMP16 PDFs, both with their default
     $\alpha_s(m_Z)=0.1149$  and with the common value $\alpha_s(m_Z)=0.118$.
\label{fig:lumi-31-vs-abmp}}
\end{center}
\end{figure}
%%%%%%%%%%%%%%%%%%%%%%%%%%%%%%%%%%%%%%%%%%%%%%%%%%%%%%%%%%%%%%%%%%%%%%


%%%%%%%%%%%%%%%%%%%%%%%%%%%%%%%%%%%%%%%%%%%%%%%%%%%%%%%%%%%%%%%%%%%%%
\begin{figure}[t]
\begin{center}
  \includegraphics[scale=0.35]{images/plots/q2_31-fc-vs-pc.pdf}
  \includegraphics[scale=0.35]{images/plots/qq_31-fc-vs-pc.pdf}
  \includegraphics[scale=0.35]{images/plots/gg_31-fc-vs-pc.pdf}
   \includegraphics[scale=0.35]{images/plots/qg_31-fc-vs-pc.pdf}
   \caption{\small Same as Fig.~\ref{fig:lumi-31-vs-30}, now
     comparing  NNPDF3.1 to its modified version
   with perturbative charm.
\label{fig:lumi-31-fc-vs-pc}}
\end{center}
\end{figure}
%%%%%%%%%%%%%%%%%%%%%%%%%%%%%%%%%%%%%%%%%%%%%%%%%%%%%%%%%%%%%%%%%%%%%%

Finally, in order to further emphasize the phenomenological impact of
the new NNPDF3.1 methodology, we compare NNPDF3.1 PDFs to the modified
version in which the charm PDF is perturbatively generated, already
discussed in Sects.~\ref{sec:results-mc},~\ref{sec:phenocharm}.
Results are shown in Fig.~\ref{fig:lumi-31-fc-vs-pc}.
On the one hand, we confirm that despite having one more parametrized
PDF, uncertainties are not increased. On the other hand, the effect on
central values is moderate but non-negligible. For
 gluon-gluon and quark-gluon luminosities, differences are always below
the one-sigma level, and typically rather
less.
%
For the quark-quark channel, results do not depend on the charm
treatment
for $M_X\gsim 200$ GeV, but for smaller invariant masses
perturbatively generated charm 
leads to a larger PDF luminosity than the best-fit parametrized charm.
%
For the quark-antiquark luminosity, we find a similar pattern at small
$M_X$, but  also 
some differences at medium and large $M_X$.




\clearpage

\subsection{$W$ and $Z$ production at the LHC 13 TeV}
\label{sec:pheno13TeV}

We compare 
theoretical predictions based on the NNPDF3.1 set to  $W$ and $Z$
production data at $\sqrt{s}=13$ TeV from  
ATLAS~\cite{Aad:2016naf}. Similar measurements by CMS~\cite{CMS-PAS-SMP-15-004} are not 
included in this comparison as they are still preliminary.
%
We compute  fiducial cross-sections using {\tt
  FEWZ}~\cite{Gavin:2012sy} at NNLO QCD 
     accuracy, using NNPDF3.1, NNPDF3.0, CT14, MMHT14
     and ABMP16 PDFs, together with the corresponding PDF uncertainty band. All calculations 
(including ABMP16) are performed with $\alpha_s=0.118$.
     %
     Electroweak NLO corrections are  computed with
     {\tt FEWZ} for $Z$  production, and with {\tt HORACE3.2}~\cite{CarloniCalame:2007cd} for
     $W$  production.
     %
    The fiducial phase space for the $W^{\pm}$ cross-section
    measurement in ATLAS is 
    by $p_T^l\ge 25$ GeV and $|\eta_l|\le 2.5$ for the charged lepton
    transverse momentum and pseudo-rapidity, a missing energy of
    $p_T^\nu\ge 25$ GeV 
    and a $W$ transverse mass of $m_T\ge 50$ GeV.
    %
    For $Z$  production, 
    $p_T^l\ge 25$ GeV and $|\eta_l|\le 2.5$ for the charged leptons transverse
    momentum and rapidity and $66\le m_{ll} \le 116$ GeV for the dilepton
    invariant mass.
         
%%%%%%%%%%%%%%%%%%%%%%%%%%%%%%%%%%%%%%%%%%%%%%%%%%%%%%%%%%%%%%%%%%%%%
\begin{figure}[t]
\begin{center}
  \includegraphics[scale=0.88]{images/plots/WPWM_ratio_ATLAS_comparison_shaded.pdf}
  \includegraphics[scale=0.88]{images/plots/WZ_ratio_ATLAS_comparison_shaded.pdf}
  \caption{\small Comparison of the ATLAS 
    measurements of the $W^+/W^-$ ratio (left) 
    and the $W/Z$ ratio (right) at $\sqrt{s}=13$ TeV with theoretical
    predictions computed with different NNLO PDF sets.
    %
    Predictions are shown with (heavy) and without (light) NLO EW corrections 
    computed with {\tt FEWZ} and {\tt HORACE}, as described in the text.
\label{fig:LHC-13tev-incxsec-1}}
\end{center}
\end{figure}
%%%%%%%%%%%%%%%%%%%%%%%%%%%%%%%%%%%%%%%%%%%%%%%%%%%%%%%%%%%%%%%%%%%%%%

In Fig.~\ref{fig:LHC-13tev-incxsec-1} we compare the 
ATLAS~\cite{Aad:2016naf} 13~TeV
    measurements of the the $W^+/W^-$ and $W/Z$ ratios
    in the fiducial region at $\sqrt{s}=13$ TeV to theoretical
    predictions, both  with and without electroweak corrections.
    %    
We see that for both the $W^+/W^-$ ratio and the $W/Z$ ratio, all the PDF sets are in 
reasonably good agreement with the data. The uncertainty in the
theoretical prediction shown in the plot is the PDF  
uncertainty only: parametric uncertainties (in the values of $\alpha_s$ and $m_c$) and missing higher 
order QCD uncertainties are not included.
Interestingly, electroweak corrections shift the
theory predictions by around 0.5\%, and for all PDF sets they 
improve the agreement with the ATLAS 
measurements.
%
NNPDF3.1 results have  smaller PDF
uncertainties than NNPDF3.0, and  are in better agreement with the
ATLAS data.

The corresponding absolute $W^+$, $W^-$ and $Z$ cross-sections are
shown in Fig.~\ref{fig:LHC-13tev-incxsec-2},
    normalized in each case to the experimental central value. 
Again, predictions are generally 
    in agreement with the data, with the possible exception of 
    ABMP16 for  $Z$ production. In comparison to cross-section
    ratios, the effect of electroweak corrections on absolute
    cross-sections, around 1\% for $W^+$ and $W^-$ and around 0.5\%
for  $Z$ production, is rather less significant on the scale of
the uncertainties involved, and it does not necessarily lead to
improved agreement.
    %
Comparing  NNPDF3.1 to NNPDF3.0 we see again considerably reduced
uncertainties and improved agreement of the prediction with data. 
This improved agreement  is particularly marked for $Z$ production,
where 3.0 was about 5\% below the data,  
while now 3.1 agrees within uncertainties.
%


%%%%%%%%%%%%%%%%%%%%%%%%%%%%%%%%%%%%%%%%%%%%%%%%%%%%%%%%%%%%%%%%%%%%%
\begin{figure}[t]
\begin{center}
  \includegraphics[scale=1.0]{images/plots/Absolute_ATLAS_shaded.pdf}
  \caption{\small Same as Fig.~\ref{fig:LHC-13tev-incxsec-1}, now 
    for the absolute $W^+$, $W^-$ and $Z$ cross-sections.
    All predictions are normalized to the experimental central value.
\label{fig:LHC-13tev-incxsec-2}}
\end{center}
\end{figure}
%%%%%%%%%%%%%%%%%%%%%%%%%%%%%%%%%%%%%%%%%%%%%%%%%%%%%%%%%%%%%%%%%%%%%%

%%%%%%%%%%%%%%%%%%%%%%%%%%%%%%%%%%%%%%%%%%%%%%%%%%%%%%%%%%%%%%%%%%%%%
\begin{figure}[t]
\begin{center}
  \includegraphics[scale=1.0]{images/plots/ggH_ratio.pdf}
  \includegraphics[scale=1.0]{images/plots/ttH_ratio.pdf}
   \includegraphics[scale=1.0]{images/plots/WH_ratio.pdf}
  \includegraphics[scale=1.0]{images/plots/ZH_ratio.pdf}
  \caption{\small PDF dependence of the 
Higgs  production cross-sections at the LHC
    $13$ TeV for gluon fusion, $t\bar{t}$
    associated production, and $VH$ associated production.
    All results are shown as ratios to the central NNPDF3.1 result.
    Only PDF uncertainties are shown.
  \label{fig:ggH}
  }
\end{center}
\end{figure}
%%%%%%%%%%%%%%%%%%%%%%%%%%%%%%%%%%%%%%%%%%%%%%%%%%%%%%%%%%%%%%%%%%%%%%

%%%%%%%%%%%%%%%%%%%%%%%%%%%%%%%%%%%%%%%%%%%%%%%%%%%%%%%%%%%%%%%%%%%%%
\begin{figure}[t]
\begin{center}
  \includegraphics[scale=1.0]{images/plots/vbf_ratio.pdf}
  \includegraphics[scale=1.0]{images/plots/HH_ratio.pdf}
  \caption{\small Same as Fig.~\ref{fig:ggH} for single Higgs production
    in vector boson fusion (left) and  double Higgs production
    in gluon fusion (right).
  \label{fig:ggH2}
  }
\end{center}
\end{figure}
%%%%%%%%%%%%%%%%%%%%%%%%%%%%%%%%%%%%%%%%%%%%%%%%%%%%%%%%%%%%%%%%%%%%%%


\subsection{Higgs production}

We finally study the PDF dependence of predictions for 
inclusive Higgs production at LHC $13$ TeV, and for
Higgs pair production, which
could also be within reach of the LHC in the
near future~\cite{Bishara:2016kjn,Behr:2015oqq}.
We study single Higgs in  gluon fusion, associated production with gauge bosons and top
pairs and vector boson fusion, and double Higgs production in gluon
fusion. In each case, we show predictions normalized to the NNPDF3.1
result, and only show PDF uncertainties. All calculations 
(including ABMP16) are performed with $\alpha_s=0.118$.

The settings are the following. 
For  gluon fusion we perform the calculation at N$^3$LO using
{\tt  ggHiggs}~\cite{Ball:2013bra,Bonvini:2014jma,Bonvini:2016frm}.
Renormalization and factorization scales are set to $\mu_F=\mu_R=m_h/2$
and the computation is performed using rescaled effective theory.
For associate production with a $t\bar{t}$ pair we use
{\tt MadGraph5\_aMC@NLO}~\cite{Alwall:2014hca}, with default
factorization and renormalization scales $\mu_R=\mu_F=H_T/2$, where $H_T$ is 
the sum of the transverse masses.
For associate production with an electroweak gauge boson we use
{\tt vh@nnlo} code~\cite{Brein:2012ne} at NNLO with default scale settings.
For  vector boson fusion we perform the calculation at N$^3$LO using
{\tt proVBFH}~\cite{Dreyer:2016oyx,Cacciari:2015jma}
with the default scale settings. Finally, for double Higgs production
at the FCC~100~TeV the calculation is performed using
{\tt MadGraph\_aMC@NLO}. 

Results are shown in Figs.~\ref{fig:ggH}-\ref{fig:ggH2}.
For gluon fusion and $t\bar{t}h$, which are both driven by the gluon PDF, the
former for $x\sim10^{-2}$, and the latter for large $x$, 
results from the various  PDF sets agree within uncertainties;
NNPDF3.0 and NNPDF3.1 are also in good agreement, with the new prediction
exhibiting reduced uncertainties.
The spread of results is somewhat larger for associate production with
gauge bosons.
The NNPDF3.1 prediction is about 3\% higher than the NNPDF3.0 one,
with uncertainties reduced by a factor 2, so the two cross-sections
barely agree within uncertainties. Also, of the three PDF sets
entering the PDF4LHC15 combination, 
NNPDF3.0 gave the smallest cross-section, but NNPDF3.1 now gives the highest
one: $VH$ production is driven by the quark-antiquark luminosity,
and this enhancement for $M_X\simeq 200$ GeV between 3.0 and 3.1
could indeed be observed already in Fig~\ref{fig:lumi-31-vs-30}.
For VBF we also find that the NNPDF3.1 result is larger, by
about 2\%,  than
the NNPDF3.0 one, with smaller
uncertainties, and it is in better  agreement with other PDF sets.
Finally, for double Higgs production in gluon fusion
the central value  with  NNPDF3.1 increases slightly but is otherwise consistent
with the NNPDF3.0 prediction, and
here there  is also good agreement for all the PDF sets.



